\chapter*{General Conclusion}
\addcontentsline{toc}{chapter}{General Conclusion}

This Master Diploma report presented the scientific study of Agrio, a Digital Twin-based intelligent irrigation platform for precision agriculture. The work focused on the background, state of the art, and synthesis required to understand how Digital Twins can support smart irrigation and agricultural water management.

The background chapters showed that smart irrigation is motivated by water scarcity, climate variability, and the need to improve irrigation decisions. Precision agriculture, IoT, wireless communication, artificial intelligence, simulation, and human-in-the-loop decision support form the main technological foundation of this field. Digital Twin technology extends classical monitoring by connecting the physical state of an agricultural system with a virtual representation able to support prediction, simulation, recommendation, and feedback.

The state-of-the-art review showed that Digital Twins are increasingly studied in agriculture but remain difficult to deploy due to data heterogeneity, biological complexity, cost, connectivity, calibration, and limited long-term validation. The literature also shows that many systems focus on monitoring or decision support but do not always implement a complete sensing-to-actuation loop.

The synthesis and comparative analysis positioned Agrio as an integrated project that combines concepts from smart irrigation, agricultural Digital Twins, IoT, machine learning, simulation, human validation, and actuation. This positioning does not claim that Agrio is a fully mature field-deployed Digital Twin. Instead, it identifies the scientific gap that the engineering prototype addresses: building a coherent bridge between research concepts and an operational cyber-physical irrigation workflow.

Future scientific work should extend the review with new field-validated systems, compare additional machine learning and optimization approaches, study agronomic calibration in greater depth, and evaluate Digital Twin maturity using long-term agricultural deployments.
